\documentclass[11pt,a4paper]{article}

% ---------- packages ----------
\usepackage[T1]{fontenc}
\usepackage[utf8]{inputenc}
\usepackage[margin=2.3cm]{geometry}
\usepackage[table]{xcolor}
\usepackage{booktabs}
\usepackage{array}
\usepackage{tabularx}
\usepackage{enumitem}
\usepackage{titlesec}
\usepackage{fancyhdr}
\usepackage{graphicx}
\usepackage[hidelinks]{hyperref}

% ---------- colours ----------
\definecolor{navy}{HTML}{000000}
\definecolor{accent}{HTML}{000000}
\definecolor{lightblue}{HTML}{F2F2F2}
\definecolor{lightgreen}{HTML}{E8E8E8}
\definecolor{amber}{HTML}{ECECEC}
\definecolor{greyln}{HTML}{808080}

% ---------- section styling ----------
\titleformat{\section}{\Large\bfseries\color{navy}}{\thesection.}{0.6em}{}[{\color{accent}\titlerule[1pt]}]
\titleformat{\subsection}{\large\bfseries\color{black!80}}{\thesubsection}{0.6em}{}
\titlespacing*{\section}{0pt}{14pt}{6pt}

% ---------- header / footer ----------
\pagestyle{fancy}\fancyhf{}
\renewcommand{\headrulewidth}{0.3pt}
\lhead{\small\color{navy}AMR-X Prototype V1}
\rhead{\small\color{navy}Drive-Wheel Suspension Study}
\cfoot{\small\color{black!55}Page \thepage}

\newcolumntype{Y}{>{\raggedright\arraybackslash}X}
\newcommand{\prio}[1]{\colorbox{amber}{\parbox{\dimexpr\linewidth-2\fboxsep}{#1}}}

\setlist[itemize]{leftmargin=1.3em,itemsep=2pt,topsep=3pt}

% ---------- title ----------
\begin{document}
\begin{titlepage}
\thispagestyle{empty}
\vspace*{2.2cm}
\noindent\rule{\linewidth}{1.2pt}\\[10pt]
{\small\bfseries\color{black} AMR-X PROTOTYPE V1 \quad$\cdot$\quad MECHANICAL ENGINEERING}\\[14pt]
{\color{black}\Huge\bfseries Drive-Wheel Suspension}\\[8pt]
{\color{black}\Large Options, Trade-offs, Recommendation and Procurement}\\[12pt]
\noindent\rule{\linewidth}{0.5pt}\\[10pt]
{\normalsize 2-drive / 4-caster differential AMR \quad$\cdot$\quad 130\,kg design load \quad$\cdot$\quad 24\,V \quad$\cdot$\quad 0.8--1.5\,m/s}
\vfill
\noindent\rule{\linewidth}{0.5pt}\\[4pt]
{\footnotesize Engineering design report \hfill Revision 1}
\end{titlepage}

% =====================================================================
\section*{List of Abbreviations}
\begin{center}\small
\renewcommand{\arraystretch}{1.25}
\begin{tabularx}{\linewidth}{@{}p{2.3cm}Y@{}}
\toprule
\textbf{Abbreviation} & \textbf{Meaning} \\
\midrule
AMR & Autonomous Mobile Robot \\
AMR-X & Designation of this prototype platform (the robot studied in this report) \\
AGV & Automated Guided Vehicle \\
\rowcolor{lightblue} PU & Polyurethane (caster/wheel tyre material) \\
RFQ & Request for Quotation \\
\rowcolor{lightblue} SIL & Safety Integrity Level (functional-safety classification, IEC\,61508) \\
PLd & Performance Level d (functional-safety rating, ISO\,13849) \\
\rowcolor{lightblue} CoG & Centre of Gravity \\
EC & Electronically Commutated --- brushless motor family (e.g.\ Maxon EC\,45) \\
\rowcolor{lightblue} GP & Gearhead, Planetary --- gearbox family (e.g.\ Maxon GP\,42C) \\
USPTO & United States Patent and Trademark Office \\
\rowcolor{lightblue} IEEE & Institute of Electrical and Electronics Engineers \\
IROS & International Conference on Intelligent Robots and Systems \\
\midrule
\multicolumn{2}{@{}l}{\textit{Units}} \\
V & Volt \\
\rowcolor{lightblue} W & Watt \\
N & Newton \\
\rowcolor{lightblue} Nm / mNm & Newton-metre / milli-newton-metre \\
kg & Kilogram \\
\rowcolor{lightblue} mm & Millimetre \\
m/s & Metres per second \\
\bottomrule
\end{tabularx}
\end{center}

% =====================================================================
\section{Why this question exists}
The layout uses \textbf{six ground contacts}: two central drive wheels plus four corner casters.
On a rigid frame that is one contact too many. Any three points define a plane, so a rigid
six-contact base behaves like a four-legged table on an uneven floor: it rocks, and whichever
contacts lift stop carrying load.

For a differential drive this is not cosmetic. Steering depends on \emph{both} drive wheels being
firmly pressed to the floor. If one unloads on an expansion-joint seam or a dock-plate lip, the wheel
slips, the robot yaws unpredictably, odometry drifts, and — with our payload centre of gravity up at
$\approx$398\,mm — the rocking is amplified. The real question is therefore not ``do we want a fancy
suspension'' but \textbf{how do we guarantee all six contacts stay loaded on a real warehouse floor.}

\textbf{Design target.} Production AMRs in this class (e.g.\ MiR) publish a traversable floor-gap
tolerance of $\approx$30\,mm at $\approx$27\,mm ground clearance~\cite{mir}. That confirms the target:
small, repeated unevenness (joints, ramps, thresholds), \emph{not} rough terrain.

% =====================================================================
\section{Candidate mechanisms}
Six approaches span the spectrum from doing nothing to full rover suspension.

\begin{center}\small
\renewcommand{\arraystretch}{1.25}
\begin{tabularx}{\linewidth}{@{}p{3.3cm}Y c c c@{}}
\toprule
\textbf{Option} & \textbf{How it works} & \textbf{Traction} & \textbf{Complexity} & \textbf{Fit} \\
\midrule
A. Rigid & all wheels bolted hard to frame & Poor & Lowest & Flat floor only \\
B. Sprung casters & springs on the corner casters & Fair & Low & Light robots \\
\rowcolor{lightgreen} C. Sprung drive units & each drive wheel on a pre-loaded spring/slide ($\sim$15--25\,mm travel) & Good & Low--Med & \textbf{Our case} \\
\rowcolor{lightgreen} D. Transverse rocker & both drive units on one central pivot beam (see-saw) & Good & Low & \textbf{Our case} \\
E. Independent arms & each drive wheel on its own sprung swing-arm & Very good & Med--High & Uneven floors \\
F. Rocker-bogie & linked rockers + body differential (Mars rover) & Excellent & High & Rough terrain (overkill) \\
\bottomrule
\end{tabularx}
\end{center}

\noindent\textbf{The two that fit us are C and D.}
\textbf{C} (sprung drive units) directly guarantees a constant downward force on each drive wheel, so
traction is maintained over seams and ramps and the payload is isolated from jolts. \textbf{D}
(transverse rocker) is the simplest passive way to keep \emph{both} drive wheels loaded: when one
rises, the beam rocks and the other is pushed down by the same amount. In practice the two are often
combined --- a sprung drive pod mounted on a rocker.

% =====================================================================
\section{Real-robot precedents}
\textbf{Option C --- sprung drive wheels.}
The \textbf{iRobot Roomba} is the textbook case: each differential drive wheel is a self-contained
spring-loaded module (motor + gearbox + spring) that compresses and extends independently to keep
floor contact; the design rationale is exactly ours --- losing drive-wheel contact loses steering
accuracy~\cite{ifixit,sevarg}. At our scale, AMR patent \textbf{US~12,128,554} places the drive units
on coil springs and linear bearings with a $\approx$25\,mm stroke and couples both wheels (``parallel''
suspension) to keep the track width constant and stay narrow-aisle compact~\cite{patentC}.
\textbf{MiR} AMRs use a central differential drive whose published gap/clearance figures only make
sense with loaded/sprung drive wheels~\cite{mir}.

\textbf{Option D --- rocker / equalizer.}
The Mars-rover rocker-bogie is the most-illustrated example of the pivot principle: if one wheel goes
up, the linked wheel is forced down by the same amount~\cite{planetary}. AMR driving-module patents use
the same single-pivot ``rocker'' idea in a compact drive unit~\cite{patentD}. We need only the
\emph{single central rocker}, not the full bogie.

% =====================================================================
\section{Recommendation: Solution C (sprung drive units)}
\textbf{We recommend Solution~C} --- placing each drive wheel on its own pre-loaded spring suspension.
It is the mechanism the field converges on for indoor differential AMRs because it attacks our exact
failure mode (a drive wheel unloading) at the source, isolates the electronics and payload from floor
shock, and --- unlike a bare rocker --- gives a known, repeatable wheel load that can be tuned to our
65\,kg-per-wheel share. It is also the only one of the two available as a \emph{finished, bolt-on
product} (Section~6), which de-risks and accelerates a prototype.

\subsection{The trade-off you must accept}
\prio{\textbf{Buying Solution C as a ready module means replacing the Maxon EC\,45 flat + GP\,42C.}
The off-the-shelf ``AGV suspension drive wheel'' is sold as a sealed assembly --- wheel + gearbox +
\emph{motor} + encoder + spring (sometimes brake and controller) --- all dimensioned around \emph{its own}
motor. There is no slot to insert the EC\,45. So adopting a ready C-module retires our discrete
motor/gearbox selection. The upside: these modules use 400\,W-class motors, which \emph{exceed} the
2$\times$70\,W Maxon set-up --- so torque and power go \emph{up}, not down (see Section~5--6).}

\medskip
\noindent If keeping the EC\,45 + GP\,42C is mandatory, Solution~C must instead be \emph{fabricated}
(a sprung vertical linear slide per drive unit, built from catalogue linear guides + springs), or
Solution~D (a fabricated rocker) is used --- both keep the Maxon and all dimensions, but neither is a
single purchasable kit.

% =====================================================================
\section{Specifications to preserve or improve}
Any replacement module must \textbf{meet or beat} the current drivetrain. Targets derived from the
sizing study:

\begin{center}\small
\renewcommand{\arraystretch}{1.25}
\begin{tabularx}{\linewidth}{@{}p{5.2cm}p{4.3cm}Y@{}}
\toprule
\textbf{Parameter} & \textbf{Current (EC\,45 + GP\,42C)} & \textbf{Requirement for a kit} \\
\midrule
Supply voltage & 24\,V & 24\,V (match battery) \\
Max speed & 1.5\,m/s (33:1 gives 2.85 top) & $\geq$1.5\,m/s \\
Continuous tractive force (pair) & $\approx$69\,N & $\geq$69\,N; $\geq$100\,N desired for ramps \\
Continuous torque (per motor) & 130\,mNm~\cite{maxon} & meet/exceed at the wheel \\
Wheel diameter & \O180\,mm & \O180\,mm to preserve geometry \\
Per-wheel load (spring tuning) & $\approx$65\,kg & spring rated/tuned for $\sim$65\,kg/wheel \\
Encoder & yes (Hall / E6B2) & integrated encoder \\
\bottomrule
\end{tabularx}
\end{center}

\noindent A 400\,W-class integrated drive wheel clears the speed, torque and power targets with margin
(``better''); the two items that need active confirmation are \textbf{24\,V} availability and a
\textbf{spring rate tuned to our light 65\,kg/wheel} (most modules are built for far heavier robots, so
an untuned spring can be too stiff to deflect).

% =====================================================================
\section{Ready-to-buy kits (Solution C) suited to our use case}
The following provide drive + suspension in one unit and are appropriate (with the noted configuration)
for a 130\,kg indoor AMR. All are customisable with low minimum order quantity.

\begin{center}\small
\renewcommand{\arraystretch}{1.3}
\begin{tabularx}{\linewidth}{@{}p{3.0cm}Y p{2.6cm}@{}}
\toprule
\textbf{Product} & \textbf{Why it fits us / what to configure} & \textbf{Suspension} \\
\midrule
\rowcolor{lightblue}
ZHLUN Robotic Drive Unit (electric differential wheel)~\cite{zhlun} &
Bundles motor, encoder, reducer, controller and PU wheel; two-stage damping (coil spring + PU block).
Request 24\,V, \O180 wheel, $\geq$1.5\,m/s, spring tuned to 65\,kg/wheel. & Spring + PU block (built in) \\
TZBOT AGV Suspension Drive Wheel~\cite{tzbot} &
Dedicated suspension differential unit; select the \emph{smallest} model (the TZCS heavy-load variants
are 4-tonne class --- avoid). Confirm 24\,V and wheel \O. & Spring suspension (built in) \\
\rowcolor{lightblue}
BeUDMKE Shock-Absorbing Drive Wheel~\cite{dmk} &
One-stop wheel + motor + driver + gearbox with shock absorption; customisable voltage, wheel \O\ and
ratio. & Shock-absorbing (built in) \\
ez-Wheel SWD\,125~\cite{ezwheel} &
24\,V integrated \emph{safety} wheel drive purpose-built for AMRs (SIL2/PLd), encoder included ---
strong on functional safety. Note: \O125 wheel (changes ride height) and confirm whether compliance is
built in or added. & Integrated drive (verify) \\
Bosch Rexroth ROKIT Motor~\cite{rokit} &
Premium plug-and-drive wheel (wheel, gearbox, brake, motor, encoder, controller). Pallet-scale ---
likely oversized; rigid, so pair with compliance. Reference-grade fallback. & Rigid (add compliance) \\
\bottomrule
\end{tabularx}
\end{center}

\noindent\textbf{Primary picks for us:} the \textbf{ZHLUN} and \textbf{BeUDMKE} suspension drive wheels,
because they explicitly include the spring/PU compliance (true Solution~C), are customisable to 24\,V /
\O180 / $\geq$1.5\,m/s, and carry far more than 65\,kg so the motor specs land \emph{above} ours. The
single open point on both is spring tuning for our light per-wheel load --- specify it in the RFQ.

% =====================================================================
\section{Priority of this modification}
\prio{\textbf{Priority: MEDIUM now, rising to HIGH before any real-floor deployment.}}

\begin{itemize}
\item \textbf{Not a bring-up blocker.} On a flat lab/test floor a rigid mount runs; the robot can be
powered up, driven and have its electronics, navigation and safety chain validated \emph{before} the
suspension is fitted. So it does not gate first motion.
\item \textbf{Becomes critical for deployment.} On any production floor with expansion joints,
dock-plates or ramps, the over-constrained six-contact base \emph{will} rock and intermittently unload a
drive wheel --- degrading traction, odometry and stability (worsened by the high CoG). Before field
trials this must be resolved.
\item \textbf{Ranking against other open items.} It sits \emph{below} the two hard geometric conflicts
that decide whether the chassis can be built at all --- the battery pack aspect ratio vs.\ the motors,
and the caster mount height vs.\ the 40\,mm ground clearance --- and below freezing the drivetrain
torque (the 33:1 $\rightarrow$ 51:1 ratio question). Resolve those first; schedule the suspension
decision immediately after, because choosing a ready C-module \emph{is} the drivetrain decision (it
replaces the motor) and therefore should not be left late.
\end{itemize}

\noindent\textbf{Recommended action:} issue an RFQ to ZHLUN and BeUDMKE for their smallest 24\,V
suspension differential drive wheel (\O180, $\geq$1.5\,m/s, encoder, spring tuned to $\sim$65\,kg/wheel),
and in parallel keep the fabricated rocker (Solution~D around the existing Maxon) as the fallback that
changes no specification and no dimension.

% =====================================================================
\begin{thebibliography}{99}
\small
\bibitem{maxon} maxon group. \textit{EC\,45 flat \O45\,mm brushless, 70\,W} — product/datasheet
(max.\ continuous torque 130\,mNm, up to 85\% efficiency, 24\,V winding).
\url{https://www.maxongroup.com/}
\bibitem{ifixit} iFixit. \textit{Roomba 650 Wheel Module Disassembly} (identifies the suspension pivot
pin). \url{https://www.ifixit.com/Guide/Roomba+650+wheel+module+Disassembly/97733}
\bibitem{sevarg} Sevarg. \textit{Roomba i7 Teardown} (drive module = motor + geartrain + suspension
spring, with large travel). \url{https://www.sevarg.net/2018/12/30/roomba-i7-teardown-with-the-waving-cat/}
\bibitem{patentC} USPTO. \textit{Autonomous mobile robot with a single modular platform},
US~12,128,554 (coil-spring drive units on linear bearings, $\approx$25\,mm stroke, parallel suspension).
\bibitem{patentD} USPTO. \textit{Driving module of autonomous mobile robot}, US~12,194,800
(single-pivot rocker drive unit). 
\bibitem{planetary} The Planetary Society. \textit{How the Mars rover rocker-bogie suspension works}
(``if the left front wheel goes up, the right front wheel is forced down by the same amount'').
\url{https://www.planetary.org/}
\bibitem{mir} Mobile Industrial Robots. \textit{MiR600 specifications} (ground clearance $\approx$27\,mm;
traversable floor gap up to $\approx$30\,mm). \url{https://www.mobile-industrial-robots.com/}
\bibitem{zhlun} ZHLUN / Guangzhou Wisdom Wheel. \textit{Robotic Drive Units — Electric Differential
Wheel} (two-stage spring + PU damping). \url{https://www.agv-wheel.com/}
\bibitem{tzbot} TZBOT. \textit{AGV Drive Wheel — Suspension (differential)}.
\url{https://www.tzbotautomation.com/agv-drive-wheels/differential-drive-wheel/agv-drive-wheel-suspension.html}
\bibitem{dmk} BeUDMKE (dmkmotor). \textit{AGV Drive Wheel — shock-absorbing drive wheel}.
\url{https://www.dmkmotor.com/agv-drive-wheel/}
\bibitem{ezwheel} ez-Wheel. \textit{SWD\,125 Safety Wheel Drive} (24\,V, SIL2/PLd, for AGV/AMR), via
RobotShop. \url{https://www.robotshop.com/products/swd-125-safety-wheel-drive-agv-amr}
\bibitem{rokit} Bosch Rexroth. \textit{ROKIT Motor — all-in-one wheel drive for AMRs}.
\url{https://www.boschrexroth.com/}
\bibitem{ieee} IEEE/IROS. \textit{Flexible (multilayered) suspension mechanism for a differential-drive
mobile robot} (passive spring+damper, self-levelling; for tall, tip-prone robots).
\end{thebibliography}

\end{document}
